\section{The limiting model and the perturbation contradiction}
\label{sec:perturbation}

We now show that the ordinary program constructed in
Proposition~\ref{prop:scaling-construction} is finite.  An infinite program
would give, by Proposition~\ref{prop:terminal-limit}, a decreasing family of
boundaries on small modifications of one terminal pair.  We first make the
limiting ordinary adjoint nef.  We then compare it with one sufficiently
nearby generalized adjoint.  The decisive point is that this comparison can
be made without changing the Cartier indices of finitely many rational
nef adjoints.  The rational decomposition and perturbation method follows
Birkar's argument \cite[Remark~3.1, Lemma~3.2, and Section~4]{Birkar2009}.
We give the details needed for the fixed Kähler b-part and for exact
line-bundle descent on the analytic contractions.

\subsection{A nef model of the limiting ordinary pair}

Recall the data supplied by Proposition~\ref{prop:terminal-limit} under the
assumption that the scaling program is infinite.  After truncation, all
its steps are flips.  There are strongly $\Q$-factorial compact Kähler
models $\mu_i:Y_i\to X_i$, isomorphic to one another in codimension one,
and effective boundaries $\Theta_i$ such that
\begin{equation}
 Q_{i,Y_i}:=\mu_i^*L_i(\lambda_i)
       =K_{Y_i}+\Theta_i+\lambda_i H_{Y_i}
 \quad\text{is nef}.
 \label{eq:extracted-nef}
\end{equation}
Here $L_i(t)=D_i+tH_{X_i}$, $D_i=K_{X_i}+\Delta_i$, and $\lambda_i\to0$.
The nef b-part $\mathbf H$ is the fixed datum determined by the initial
Kähler form.  On a fixed first model $Y$, the transforms $\Theta_{i,Y}$
decrease coefficientwise, on a fixed finite support, to an effective real
boundary $\Theta_Y$.  The ordinary pair $(Y,\Theta_Y)$ is terminal.
The transform of $\Theta_Y$ on each $Y_i$ pushes forward to $\Delta_i$.

\begin{proposition}
\label{prop:limiting-nef}
For these data there is a finite sequence of ordinary
$(K_Y+\Theta_Y)$-flips
\[
 Y\dashrightarrow W
\]
such that $W$ is strongly $\Q$-factorial and compact Kähler, the transformed
ordinary pair $(W,\Theta_W)$ is terminal, and
$P:=K_W+\Theta_W$ is nef.  In particular, $W$ is isomorphic in codimension
one to every $Y_i$.
\end{proposition}

\begin{proof}
Run the ordinary real-boundary program from $(Y,\Theta_Y)$ using
\cite[Proposition~4.2]{HX2026}, with zero nef part.
We show inductively that every step is a flip.  Let $T$ be a current
model, so that $T$ is small to every $Y_i$, and write
\[
 P_T=K_T+\Theta_T,\qquad
 Q_{i,T}=K_T+\Theta_{i,T}+\lambda_i H_T .
\]
The trace $H_T$ exists by Lemma~\ref{lem:class-trace}: along these
strongly $\Q$-factorial real-boundary flips we use the canonical
Bott--Chern transport of \cite[Proposition~4.2]{HX2026}.  Its pullback
comparison has an actual exceptional-divisor correction, so the same
trace relation propagates with the fixed b-data.
Together with the effective transformed boundary, the actual trace
relation supplies the generalized structure boundary for $Q_{i,T}$;
no generalized klt condition is required for the small-model comparison
below. The finite boundary support and
$\lambda_i\to0$ give $Q_{i,T}\to P_T$ in $\BC(T)$.

Fix a $P_T$-negative extremal ray and its contraction.  For all sufficiently
large $i$ this same ray is $Q_{i,T}$-negative.  Take a common smooth
Kähler resolution $p:U\to T$, $q:U\to Y_i$.  The boundaries are transforms
of one another and the b-part is unchanged.  Thus
Lemma~\ref{lem:small-comparison}, together with
\eqref{eq:extracted-nef}, gives an actual divisor $E_i$ such that
\begin{equation}
 p^*Q_{i,T}=q^*Q_{i,Y_i}+[E_i],
 \qquad E_i\ge0,
 \qquad E_i\text{ is exceptional over both models}.
 \label{eq:limit-comparison}
\end{equation}

Suppose that the contraction were divisorial.  Choose a general point
of an exceptional prime outside the codimension-at-least-two image of
$\Exc(p)$.  A positive-dimensional projective fibre through this point
contains an irreducible curve $C$.  Its strict transform $\widetilde C$
on $U$ is not contained in $\Supp(E_i)$.  Consequently
\[
 Q_{i,T}\cdot C
   =q^*Q_{i,Y_i}\cdot\widetilde C+E_i\cdot\widetilde C\ge0,
\]
contrary to the choice of the ray.  For a fibre-type contraction the
same argument uses a general point of $T$ instead.  It also rules out
ending the program with a Mori fibre space.

Therefore every step is an ordinary flip.  Terminality persists by the
ordinary discrepancy comparison.  These flips have the exact relative
properties needed for the auxiliary termination theorem.  The negative
contraction is projective by \cite[Proposition~4.2]{HX2026}; the positive
one is proper birational between rational-singularity spaces, hence
strongly Moishezon by \cite[Lemma~2.39]{HX2026}, and projective by
\cite[Lemma~2.40]{HX2026}.  For a step
$T\xrightarrow{f}Z\xleftarrow{f^+}T^+$, the signed ordinary adjoints
$-P_T$ and $P_{T^+}$ are relatively Kähler real line bundles.
To obtain relative ampleness in this intrinsic line-bundle setting,
write either signed adjoint $L$ in a finite real span of integral line
bundles, and choose a base class $\gamma$ such that
$c_1(L)+h^*\gamma$ is Kähler, where $h=f$ or $f^+$.
Openness of the Kähler cone gives a rational simplex around the
coefficient vector of $L$ whose rational vertices $L_\nu$ still have
$c_1(L_\nu)+h^*\gamma$ Kähler.  Lemma~\ref{lem:relative-positivity}
makes every $L_\nu$ relatively ample, and their positive convex
combination $L$ is therefore relatively ample as a real line bundle.  Thus
Proposition~\ref{prop:terminal-real-four} makes this sequence finite.
Its endpoint has nef ordinary adjoint, as asserted.
\end{proof}

\subsection{Rational nef boundaries and a positive intersection gap}

The following argument gives a rational description near a nef ordinary
adjoint, even though the limiting boundary has real coefficients.
Only the ordinary boundary coefficients vary in this argument; the
transcendental b-part plays no role.

\begin{lemma}
\label{lem:rational-nef-polytope}
Let $W$ be a globally $\Q$-factorial compact Kähler fourfold with $K_W$
a $\Q$-line bundle.  Let $B\ge0$ be a real boundary such that $(W,B)$ is
klt and $K_W+B$ is nef.  There are effective rational klt boundaries
$B^1,\ldots,B^r$ and positive real numbers $u_1,\ldots,u_r$ such that
\[
 \sum_{j=1}^r u_j=1,\qquad
 B=\sum_{j=1}^r u_jB^j,\qquad
 K_W+B^j\ \text{is nef for every }j.
\]
The boundaries $B^j$ can be chosen on the positive support of $B$.
\end{lemma}

\begin{proof}
If $B=0$, take $r=1$ and $B^1=0$.  Otherwise work in the finite-dimensional
coefficient space on $\Supp(B)$.  A sufficiently small simplex with
rational vertices $C^0,\ldots,C^d$ contains $B$ in its interior and
consists of effective klt boundaries.  Indeed, all coefficients of $B$
on this support are positive, and the klt inequalities are open on one
log resolution.  Put $D_k=K_W+C^k$.

Choose a rational number $\delta>0$ smaller than all the barycentric
coordinates of $B$.  Let $\mathcal P$ be the smaller rational polytope
in the simplex defined by requiring every barycentric coordinate to
be at least $\delta$.  We claim that its nef locus
\[
 \mathcal N=\{C\in\mathcal P:K_W+C\text{ is nef}\}
\]
is a rational polytope.

Consider an extremal ray $R$ that is negative for $K_W+C$ for some
$C\in\mathcal P$.  At least one vertex adjoint, say $D_a$, is negative
on $R$.  Among all irreducible rational curves whose classes span $R$,
choose $\Gamma$ minimizing $-D_a\cdot\Gamma$.  Such curves exist by the
cone theorem \cite[Theorem~1.3]{HX2026}, and a minimum exists because
a Cartier multiple of $D_a$ puts their positive degrees in a discrete
subset of $\R_{>0}$.

This choice minimizes the numerical scale of the curve class on $R$.
For any other vertex $D_k$ negative on $R$, the cone theorem supplies
a rational curve $\Gamma_k$ on $R$ with
$0<-D_k\cdot\Gamma_k\le8$.  Since $[\Gamma]$ is a positive multiple
of $[\Gamma_k]$ with multiplier at most one, it follows that
$-D_k\cdot\Gamma\le8$ as well.  Vertices nonnegative on $R$ need no
estimate.  Thus all the numbers $d_k=D_k\cdot\Gamma$ satisfy $d_k\ge-8$.

Write a member of $\mathcal P$ that is negative on $R$ as
$\sum t_kC^k$, where $t_k\ge\delta$ and $\sum t_k=1$.  Then
$\sum t_kd_k<0$, so for each $k$,
\[
 t_kd_k<8(1-t_k).
\]
In particular,
\begin{equation}
 -8\le d_k<\frac8\delta\qquad(0\le k\le d).
 \label{eq:finite-intersection-vectors}
\end{equation}
Each $D_k$ has a fixed global Cartier index on $W$.  Hence its degrees
belong to a fixed discrete subgroup of $\R$, and
\eqref{eq:finite-intersection-vectors} permits only finitely many
vectors $(d_0,\ldots,d_d)$.  Their nonnegativity inequalities are rational.

These finitely many inequalities define exactly $\mathcal N$ inside
$\mathcal P$.  Indeed, a ray never negative on $\mathcal P$ imposes
no further condition.  For any other ray the chosen vector tests its
nonnegativity.  Finally, positivity on all extremal rays is positivity
on $\NA(W)$: intersect this cone with the hyperplane where a fixed
Kähler class has value one and use the resulting compact convex slice.
Nef-cone duality now proves the claim.  Express $B$ as a convex
combination of rational vertices of $\mathcal N$, omitting vertices
with zero weight.  These are the required $B^j$.
\end{proof}

Apply Lemma~\ref{lem:rational-nef-polytope} to $(W,\Theta_W)$ from
Proposition~\ref{prop:limiting-nef}.  Fix a decomposition
\begin{equation}
 \Theta_W=\sum_{j=1}^r u_jB^j,\qquad
 P=\sum_{j=1}^r u_jP^j,\qquad P^j=K_W+B^j,
 \label{eq:nef-decomposition}
\end{equation}
and positive integers $\ell_j$ such that $\ell_jP^j$ is an integral
line bundle, also clearing the coefficients of $B^j$.
Define
\begin{equation}
 c=\min_{1\le j\le r}\frac{u_j}{\ell_j}>0,
 \qquad
 m\in\N,\quad (m-1)c>8.
 \label{eq:gap-and-amplification}
\end{equation}
Whenever the transforms of the $P^j$ remain nef and the same $\ell_j$
remain Cartier indices, every curve $\Gamma$ satisfies
\begin{equation}
 P\cdot\Gamma>0\quad\Longrightarrow\quad P\cdot\Gamma\ge c.
 \label{eq:positive-gap}
\end{equation}
At least one nonnegative summand in \eqref{eq:nef-decomposition} is then
positive, and its degree is at least $1/\ell_j$.
The next argument preserves exactly these indices.

\subsection{A nearby program that is trivial for the limiting class}

\begin{proposition}
\label{prop:trivial-perturbation}
For the extracted data of Proposition~\ref{prop:terminal-limit} and the
model $W$ of Proposition~\ref{prop:limiting-nef}, take $i$ sufficiently
large.  There is a finite generalized $Q_{i,W}$-program
$W\dashrightarrow W_i$ consisting only of flips, with the following
properties:
\begin{enumerate}
 \item $Q_{i,W_i}$ is nef;
 \item every step is trivial for the transformed class $P$ and for
       each $P^j$ in \eqref{eq:nef-decomposition};
 \item the transformed $P^j$ are nef adjoints of ordinary rational
       klt pairs, and the fixed multiples $\ell_jP^j$ are line bundles;
 \item every $Q_{i,W_i}$-trivial curve is $P_{W_i}$-trivial.
\end{enumerate}
\end{proposition}

\begin{proof}
On $W$, set $G_i=\Theta_{i,W}-\Theta_W\ge0$.  Keep the integer $m$ in
\eqref{eq:gap-and-amplification} fixed and consider the generalized
adjoint
\begin{equation}
 A_W=K_W+\Theta_W+mG_i+m\lambda_iH_W
     =mQ_{i,W}-(m-1)P .
 \label{eq:amplified-adjoint}
\end{equation}
For large $i$ its pair is gklt.  To see this, choose one log resolution
carrying the finite boundary support and the fixed b-data.  The
crepant boundary coefficients vary continuously with $G_i$ and
$\lambda_i$; at the limit they are those of the ordinary klt pair
$(W,\Theta_W)$.  Thus they remain strictly less than one.
The boundary $\Theta_W+mG_i$ is effective.  The pair for $Q_{i,W}$
is the convex interpolation with weight $1/m$ between this pair and
$(W,\Theta_W)$, so it is gklt too.

The sum $\Theta_{i,W}+\lambda_iH_W$ is big:
$\lambda_i>0$, the trace $H_W$ is big by Corollary~\ref{cor:big-trace},
and $\Theta_{i,W}$ is effective.  The model $W$ is strongly
$\Q$-factorial.  Choose a sufficiently large Kähler scaling class and
apply \cite[Theorem~1.5]{HX2026} to obtain a finite $Q_{i,W}$-program.

All its steps are flips and its endpoint is nef.  Indeed, as long as
the current model $T$ is small to $Y_i$, the comparison
\[
 p^*Q_{i,T}=q^*Q_{i,Y_i}+[E],\qquad E\ge0
\]
holds on a common resolution, with $E$ exceptional over both models.
The curve-through-a-general-point argument in
Proposition~\ref{prop:limiting-nef} excludes a negative divisorial
contraction or a fibre-type contraction.  It applies with this fixed
$i$, without any limiting argument.

We prove the remaining assertions by induction along these flips.
Suppose that on the current model $T$ the ordinary pairs
$(T,B^j_T)$ are klt, the $P^j_T=K_T+B^j_T$ are nef with
$\ell_jP^j_T$ line bundles, and the amplified pair with class
\[
 A_T=mQ_{i,T}-(m-1)P_T
\]
is gklt.  These conditions hold initially.  A $Q_{i,T}$-negative
extremal ray is $A_T$-negative because $P_T$ is nef.  The cone theorem
for the amplified pair gives a rational curve $\Gamma$ spanning that
ray with $A_T\cdot\Gamma\ge-8$.  If $P_T\cdot\Gamma>0$, then
\eqref{eq:positive-gap} implies
\[
 mQ_{i,T}\cdot\Gamma
   =A_T\cdot\Gamma+(m-1)P_T\cdot\Gamma
   \ge-8+(m-1)c>0,
\]
a contradiction.  Hence the ray is $P_T$-trivial.  Since every $u_j$
is positive and every $P^j_T$ is nef, it is also $P^j_T$-trivial for
each $j$.

Let $f:T\to Z\leftarrow T^+:f^+$ be this flip.  For the ordinary
rational klt pair $(T,B^j_T)$, the line bundle
$M_j=\ell_jP^j_T$ is numerically trivial over $Z$, and
\[
 M_j-(K_T+B^j_T)=(\ell_j-1)P^j_T
\]
is relatively nef.  Lemma~\ref{lem:descent} gives an actual line
bundle $M_{j,Z}$ with $M_j=f^*M_{j,Z}$.
The pullback $(f^+)^*M_{j,Z}$ agrees on the common big open set with
the reflexive sheaf of $\ell_j(K_{T^+}+B^j_{T^+})$.  Normality and
reflexivity extend the agreement.  Thus the \emph{same} integer
$\ell_j$ is a Cartier index on $T^+$; no further multiple is taken.

It follows that both adjoints $P^j_T$ and $P^j_{T^+}$ pull back from
the same class on $Z$.  On a common smooth compact Kähler resolution
$p:V\to T$, $q:V\to T^+$, their pullbacks are therefore equal.
The class $p^*P^j_T$ is nef, and nef descent along $q$ in
Lemma~\ref{lem:pullback-image} shows that $P^j_{T^+}$ is nef: the target
$T^+$ is generalized klt and hence has rational singularities.
Their actual discrepancy-change divisor is zero by negativity, so
the pairs remain klt and the comparison is crepant.
In particular the same is true of their convex combination $P_T$.

Finally, on a common resolution the discrepancy-change divisor for
the amplified pair is now $m$ times that for the $Q_i$-pair, because
the contribution from $P_T$ is zero.  The latter divisor is effective
along the chosen generalized flip.  The amplified pair therefore
remains gklt.  This completes the induction, including the fixed-index
assertion and the validity of \eqref{eq:positive-gap} at the endpoint.

It remains to prove assertion (4).  Write $Q=Q_{i,W_i}$ and
$P'=P_{W_i}$.  Both are nef.  Suppose a $Q$-trivial curve has positive
$P'$-degree.  The face
\[
 \mathcal F=\{\gamma\in\NA(W_i):Q\cdot\gamma=0\}
\]
then has a compact Kähler slice on which $P'$ takes a positive value.
Maximize $P'$ on that slice and choose an extreme point of its
maximizing face.  This point is extreme in the slice of $\mathcal F$.
Since $\mathcal F$ is a face of $\NA(W_i)$, it spans an extremal ray
of the full cone.  On this ray $Q=0$ and $P'>0$, so the amplified
class $A=mQ-(m-1)P'$ is negative.  Its cone theorem gives a rational
generator $\Gamma$ with
\[
 -8\le A\cdot\Gamma
      =-(m-1)P'\cdot\Gamma
      \le-(m-1)c<-8,
\]
which is impossible.
\end{proof}

\subsection{Descent to the original flipping base}

The preceding construction has produced a nef limiting class that
vanishes on every curve killed by the nearby nef generalized class.
We now descend it to the base of one of the original flips.  This
forces the original negative adjoint to have degree zero on its
chosen ray, giving the contradiction.

\begin{proposition}
\label{prop:scaling-finite}
Let $(X,\Delta)$ satisfy the hypotheses of
Proposition~\ref{prop:scaling-construction}.  The ordinary program
constructed there, with scaling of the fixed initial Kähler b-part,
is finite.
\end{proposition}

\begin{proof}
Suppose it were infinite.  Use Proposition~\ref{prop:terminal-limit}
and choose $W$, then a sufficiently large $i$ and $W_i$, as in
Propositions~\ref{prop:limiting-nef} and
\ref{prop:trivial-perturbation}.  Let $f_i:X_i\to Z_i$ be the
contraction of the chosen $D_i$-negative ray in the original flip
tail.  The threshold class $L_i(\lambda_i)$ descends to $Z_i$.

Take a common smooth compact Kähler resolution
$p:V\to W_i$ and $q:V\to Y_i$.  Put
$s=\mu_iq$ and $r=f_is$.  The relevant maps are shown in
Figure~\ref{fig:perturbation-descent}.
\begin{figure}[htbp]
 \centering
 \begin{tikzpicture}[
   >=Latex,
   every node/.style={font=\small},
   map/.style={->,semithick}
 ]
  \node (V) at (0,2.5) {$V$};
  \node (Wi) at (-2.4,1.25) {$W_i$};
  \node (Yi) at (2.4,1.25) {$Y_i$};
  \node (Xi) at (2.4,0) {$X_i$};
  \node (Zi) at (0,-1.15) {$Z_i$};
  \draw[map] (V)--node[above left] {$p$}(Wi);
  \draw[map] (V)--node[above right] {$q$}(Yi);
  \draw[map] (Yi)--node[right] {$\mu_i$}(Xi);
  \draw[map] (Xi)--node[below right] {$f_i$}(Zi);
  \draw[map] (V)--node[left,pos=.63] {$r=f_i\mu_iq$}(Zi);
  \draw[dashed,semithick] (Wi)--node[above,pos=.28] {small}(Yi);
 \end{tikzpicture}
 \caption{The final comparison is made on $V$.
 The class $p^*P_{W_i}$ descends along $r$; no morphism from $W_i$
 to $X_i$ or to $Z_i$ is required.}
 \label{fig:perturbation-descent}
\end{figure}

The models $W_i$ and $Y_i$ are small modifications with the same
transformed boundary and b-data.  Their $Q_i$-classes are both nef.
Applying Lemma~\ref{lem:small-comparison} in both directions gives
\begin{equation}
 p^*Q_{i,W_i}=q^*Q_{i,Y_i}
            =s^*L_i(\lambda_i).
 \label{eq:equal-nef-pullbacks}
\end{equation}
Consequently this class has zero degree on every $r$-vertical curve.
For such a curve $C$, if $p(C)$ is a point then
$p^*P_{W_i}\cdot C=0$ immediately.  Otherwise the projection formula
and \eqref{eq:equal-nef-pullbacks} show that $p(C)$ is
$Q_{i,W_i}$-trivial.  Proposition~\ref{prop:trivial-perturbation}(4)
then gives $p^*P_{W_i}\cdot C=0$ as well.

The map $r$ is proper and birational.  Its target $Z_i$ is normal
compact Kähler with rational singularities by the ordinary-step
construction, and its source $V$ is smooth compact Kähler.
Its fibres are connected, and its general fibre is a point.
Thus all the hypotheses of Lemma~\ref{lem:pullback-image} apply to
$r$.  There is a class $\xi\in\BC(Z_i)$ with
\begin{equation}
 p^*P_{W_i}=r^*\xi.
 \label{eq:limiting-descent}
\end{equation}

We next form the ordinary canonical-and-boundary comparison.
Using compatible canonical data, the difference
\begin{equation}
 F=p^*(K_{W_i}+\Theta_{W_i})
          -s^*(K_{X_i}+\Delta_i)
 \label{eq:final-exceptional-divisor}
\end{equation}
is an actual real divisor on $V$, with the relative canonical terms
interpreted as in Section~\ref{sec:analytic}.
It is exceptional over $X_i$: on every prime that is not
$s$-exceptional the two boundary coefficients agree, since the
limiting boundary on $Y_i$ pushes to $\Delta_i$ and $W_i$ is small
to $Y_i$.  This constructs the divisor before any appeal to its
cohomology class.

By \eqref{eq:limiting-descent},
\[
 [F]=p^*P_{W_i}-s^*D_i
     =s^*(f_i^*\xi-D_i).
\]
Hence $F$ is numerically trivial over $X_i$.
Lemma~\ref{lem:negativity}, applied to $F$ and to $-F$, gives $F=0$.
Pullback injectivity for $s$ now yields
\[
 D_i=f_i^*\xi\quad\text{in }\BC(X_i).
\]
Its intersection with the ray contracted by $f_i$ is therefore zero,
contradicting the strict $D_i$-negativity with which that ray was chosen.
The scaling program is finite.
\end{proof}
